\documentclass[aspectratio=169, 10pt]{beamer}

% Configuración de idioma y paquetes estándar
\usepackage[spanish, es-tabla]{babel}
\usepackage{booktabs}
\usepackage{graphicx}

% Rutas de imágenes (compatible compilando desde la raíz o desde la carpeta /latex)
\graphicspath{{images/}{../images/}}

% Tema y colores corporativos
\usetheme{Madrid}
\usecolortheme{default}

\definecolor{NavyBlue}{RGB}{27, 54, 93}
\definecolor{SlateGray}{RGB}{74, 85, 104}

\setbeamercolor{palette primary}{bg=NavyBlue,fg=white}
\setbeamercolor{palette secondary}{bg=SlateGray,fg=white}
\setbeamercolor{palette tertiary}{bg=NavyBlue,fg=white}
\setbeamercolor{titlelike}{parent=palette primary}
\setbeamercolor{block title}{bg=NavyBlue,fg=white}
\setbeamercolor{block body}{bg=gray!10,fg=black}

% Ocultar símbolos de navegación inferiores
\setbeamertemplate{navigation symbols}{}

% Información de la Portada
\title[Informe Técnico JxR]{Inserción Laboral para Mujeres Anteriormente Privadas de Libertad}
\subtitle{Propuesta Metodológica Ajustada}
\author[V. Rojas \& L. Capobianco]{Vicente Rojas Murúa \and Leonardo Capobianco}
\institute[JxR - AIM Chile - UChile]{
  Fundación Juntos por la Reinserción (JxR) \par
  AIM Chile | Facultad de Ciencias Sociales, Universidad de Chile
}
\date{Octubre de 2026}

\begin{document}

% -----------------------------------------------------------------------------
% LÁMINA 1: PORTADA
% -----------------------------------------------------------------------------
\begin{frame}
  \titlepage
\end{frame}

% -----------------------------------------------------------------------------
% LÁMINA 2: CONTEXTO Y DESAFÍO
% -----------------------------------------------------------------------------
\begin{frame}{1. Contexto y Desafío}
  \begin{columns}[T]
    \begin{column}{0.48\textwidth}
      \begin{block}{Contexto}
        \begin{itemize}
          \item \textbf{162.043 personas} atendidas por Gendarmería (88,2\% hombres / 11,8\% mujeres).
          \item Unos \textbf{30.000 egresos} proyectados entre 2025 y 2030; \textbf{8,2\% son mujeres}.
          \item \textbf{85,4\%} de las mujeres son madres.
          \item Cerca del \textbf{95\%} de las personas que JxR intermedia son hombres.
        \end{itemize}
      \end{block}
    \end{column}
    \begin{column}{0.48\textwidth}
      \begin{block}{Desafío}
        \begin{itemize}
          \item \textbf{Deseabilidad social:} la corrección política en las encuestas oculta vetos implícitos.
          \item \textbf{Doble penalización:} estigma penal, estigma de género y cargas de cuidado.
          \item \textbf{Desconexión} entre la oferta de JxR y los incentivos y riesgos de cada empresa según tamaño y rubro.
        \end{itemize}
      \end{block}
    \end{column}
  \end{columns}
\end{frame}

% -----------------------------------------------------------------------------
% LÁMINA 3: PREGUNTAS Y OBJETIVOS
% -----------------------------------------------------------------------------
\begin{frame}{2. Preguntas y Objetivos}
  \begin{block}{Pregunta Central}
    \small ¿Qué factores organizacionales, condiciones habilitantes y percepciones asociadas al género ayudan a comprender cómo y por qué las empresas deciden adoptar y sostener la contratación de mujeres anteriormente privadas de libertad en la red de JxR?
  \end{block}

  \begin{block}{Objetivos Específicos}
    \small
    \begin{enumerate}
      \item \textbf{Línea base:} permanencia a 30, 60 y 90 días, por género.
      \item \textbf{Contraste hombre--mujer:} viñetas pareadas como complemento del foco en mujeres.
      \item \textbf{Condiciones de adopción:} combinaciones observadas mediante QCA exploratorio.
      \item \textbf{Procesos:} reconstruir puntos de quiebre de contrataciones sostenidas e interrumpidas.
      \item \textbf{Propuesta:} tipología B2B y ajustes a Bolsa Laboral, OTEC y CET.
      \item \textbf{Kit cuantitativo} de unos 200 casos para aplicar en 2027.
    \end{enumerate}
  \end{block}
\end{frame}

% -----------------------------------------------------------------------------
% LÁMINA 4: MÉTODO
% -----------------------------------------------------------------------------
\begin{frame}{3. Método: Estudio de Casos Múltiples en 4 Fases}
  \begin{columns}[T]
    \begin{column}{0.54\textwidth}
      \begin{block}{Fases}
        \small
        \begin{itemize}
          \item \textbf{Fase 1 (Sem 1--2):} \textit{desk research} con el catastro de unas 50 empresas de JxR y pauta validada con la fundación.
          \item \textbf{Fase 2 (Sem 3--7):} terreno con un mínimo de 10 empresas y 29 instancias.
          \item \textbf{Fase 3 (Sem 8--10):} análisis, tipología y propuesta.
          \item \textbf{Fase 4 (2027):} kit cuantitativo de unos 200 casos, según recursos.
        \end{itemize}
      \end{block}
    \end{column}
    \begin{column}{0.42\textwidth}
      \begin{block}{Muestra Cualitativa (29)}
        \small
        \begin{tabular}{l c}
          \toprule
          \textbf{Unidad} & \textbf{N} \\
          \midrule
          Casos de empresa & 18 \\
          Mujeres en reinserción & 8 \\
          Ecosistema JxR (sesión) & 1 \\
          Actores públicos & 2 \\
          \bottomrule
        \end{tabular}
      \end{block}
    \end{column}
  \end{columns}

  \begin{block}{Selección y Entrevistas}
    \small Empresas según desenlace: 3 sostuvieron, 3 interrumpieron, 3 evaluaron y no avanzaron, 1 nunca lo consideró. Las entrevistas las realiza un \textbf{equipo de profesionales}, con la participación de Vicente y Leonardo.
  \end{block}
\end{frame}

% -----------------------------------------------------------------------------
% LÁMINA 5: ANÁLISIS Y RIGOR
% -----------------------------------------------------------------------------
\begin{frame}{4. Análisis y Resguardos de Rigor}
  \begin{columns}[T]
    \begin{column}{0.48\textwidth}
      \begin{block}{QCA Exploratorio (3 condiciones)}
        \small
        \begin{enumerate}
          \item \textbf{Sponsor en la dirección.}
          \item \textbf{Compatibilidad de cuidados.}
          \item \textbf{Acompañamiento ofrecido al evaluar.}
        \end{enumerate}
        $2^3 = 8$ combinaciones posibles. Con 10 casos habrá combinaciones sin casos (diversidad limitada): se declara y se analiza solo lo observado.
      \end{block}
    \end{column}
    \begin{column}{0.48\textwidth}
      \begin{block}{Procesos y Rigor}
        \small
        \begin{itemize}
          \item Comparación de los 6 casos que contrataron, apoyada en 3 díadas empresa--trabajadora.
          \item Doble codificación de parte del material, criterio de saturación y triangulación entre empresa, trabajadora y registros de JxR.
          \item Viñetas pareadas: ejercicio proyectivo para contrastar discursos, relevante porque JxR atiende mayoritariamente a hombres.
        \end{itemize}
      \end{block}
    \end{column}
  \end{columns}
\end{frame}

% -----------------------------------------------------------------------------
% LÁMINA 6: EQUIPO Y ÉTICA
% -----------------------------------------------------------------------------
\begin{frame}{5. Equipo y Ética}
  \begin{columns}[T]
    \begin{column}{0.48\textwidth}
      \begin{block}{Equipo y Gobernanza}
        \small
        \begin{itemize}
          \item \textbf{Investigadores:} Vicente Rojas y Leonardo Capobianco (diseño, coordinación, entrevistas y análisis).
          \item \textbf{Equipo de profesionales:} realiza las entrevistas.
          \item \textbf{Panel de revisión:} Daniel Horta, Anastasia Ganderats, Valentina Piña y Rodrigo de la Riva.
          \item \textbf{Comité AIM} y \textbf{Comité Directivo de JxR}.
        \end{itemize}
      \end{block}
    \end{column}
    \begin{column}{0.48\textwidth}
      \begin{block}{Resguardos Éticos (ICC/ESOMAR)}
        \small
        \begin{itemize}
          \item \textbf{Comité de ética externo e independiente} (por confirmar); ninguna entrevista antes de su dictamen.
          \item Consentimiento diferenciado para mujeres y empresas.
          \item Díadas con confidencialidad cruzada, contención y compensación.
          \item Opt-in de contacto separado; sin prospección comercial.
        \end{itemize}
      \end{block}
    \end{column}
  \end{columns}
\end{frame}

% -----------------------------------------------------------------------------
% LÁMINA 7: CRONOGRAMA Y ENTREGABLES
% -----------------------------------------------------------------------------
\begin{frame}{6. Cronograma y Entregables}
  \begin{columns}[T]
    \begin{column}{0.50\textwidth}
      \begin{block}{Cronograma (10 semanas)}
        \footnotesize
        \begin{tabular}{l l}
          \toprule
          \textbf{Semanas} & \textbf{Etapa} \\
          \midrule
          1--2 (19--30 oct) & Desk research, pauta y ética \\
          3--7 (2 nov--4 dic) & Terreno \\
          8--9 (7--18 dic) & Análisis \\
          10 (21--24 dic) & Informe final y entrega \\
          2027 & Taller con JxR y kit \\
          \bottomrule
        \end{tabular}
      \end{block}
    \end{column}
    \begin{column}{0.46\textwidth}
      \begin{block}{5 Entregables}
        \small
        \begin{enumerate}
          \item Informe técnico y ejecutivo (cualitativo).
          \item Matriz de tipología y guía B2B de abordaje y vinculación.
          \item Ajustes a los modelos de JxR.
          \item Hoja de ruta a 90 días.
          \item Kit cuantitativo (2027).
        \end{enumerate}
      \end{block}
    \end{column}
  \end{columns}
  \vspace{0.1cm}
  \begin{center}
    \small \textbf{El estudio cualitativo se entrega la última semana de diciembre de 2026.}
  \end{center}
\end{frame}

% -----------------------------------------------------------------------------
% LÁMINA 8: SÍNTESIS Y PRÓXIMOS PASOS
% -----------------------------------------------------------------------------
\begin{frame}{7. Síntesis y Próximos Pasos}
  \begin{center}
    \Large \textbf{De la Intuición Comercial a la Evidencia Rigurosa}
  \end{center}

  \vspace{0.3cm}

  \begin{itemize}
    \item \textbf{Valor estratégico:} permite comprender cómo y por qué deciden las empresas, y priorizar modelos de colaboración, segmentos y mensajes para JxR.
    \item \textbf{Impacto social:} orienta acciones que favorezcan la contratación y la permanencia laboral de mujeres anteriormente privadas de libertad.
  \end{itemize}

  \vspace{0.3cm}

  \begin{block}{Próximos Pasos}
    Confirmar el comité de ética externo, validar la pauta de entrevista con JxR y recibir el catastro de empresas, para iniciar el terreno el 2 de noviembre.
  \end{block}
\end{frame}

\end{document}
